%!TEX root = ../dokumentation.tex

\chapter{Theorieteil}

\section{Mikrocontroller}
\subsection{Was ist ein Mikrocontroller?}
Eine \ac{MCU} ist ein kleiner Computer, bestehend aus nur einem Chip. 
Sie wurden speziell dafür entwickelt, um Aufgaben in Systemen zu übernehmen, 
ohne dafür ein komplexes Betriebssystem zu 
benötigen. 

Er enthält einen Prozessor, Arbeitsspeicher und einen programmierbaren 
Speicher (\ac{EEPROM}) sowie Ein- und Ausgabeschnittstellen.
Im Gegensatz zu Mikroprozessoren benötigt ein Mikrocontroller keine weiteren externen Komponenten.
Außerdem sind sie optimiert für die Steuerung von Sensoren oder Motoren\citeM{i}{ibmMikro}{}{}.



\subsection{Besonderheiten des ESP32}
Der ESP32 ist ein leistungsfähiger Mikrocontroller und der Nachfolger des ESP8266 von Espressif.
Er verfügt über einen Dual-Core 32-Bit Prozessor mit zwei Kernen, 
was ihn besonders schnell macht. 

WLAN gehört zu seinen Kernfunktionen. Der ESP32 kann sich mit WLAN-Netzwerken 
verbinden und zusätzlich ein eigenes WLAN-Netzwerk (Access Point) aufbauen, 
beispielsweise als kleiner Webserver. 
Zudem unterstützt er Bluetooth, was ihn für Smart-Home- und \ac{IoT}-Anwendungen ideal macht.

Der Mikrocontroller überzeugt auch durch einen sehr geringen Stromverbrauch, 
der im Deep-Sleep-Modus je nach Board unter 0,1 mA liegen kann, was ihn besonders 
für energieeffiziente \ac{IoT}-Projekte prädestiniert.

Der ESP32 ist robust und kann in einem Temperaturbereich von -40 °C bis +125 °C 
zuverlässig betrieben werden.

Entwicklungsboards sind so aufgebaut, dass sie neben solchen Modulen auch die 
notwendige Versorgung und Anschlüsse bereitstellen, 
um die vielfältigen Funktionen des ESP32 einfach nutzen zu können
\citeM{i}{digiESP.2023}{}{}.

\section{Raspberry Pi}
Was wann etc. auch aufs OS eingehen, besonders auf den neuen Network Manager

\section{Der DHT11 Sensor}
Der DHT11 Sensor ist in der Lage, Temperatur und Luftfeuchtigkeit zu messen.

Zum Messen der Temperatur hat der DHT11 ein Negative Temperature Coefficient
Thermistor. Dabei handelt es sich um einen Widerstand, welcher geringer wird je höher die 
Temperatur ist, da die Leitfähigkeit des Materials zunimmt.

Für die Luftfeuchtigkeit ist ein Kondensator in den DHT11 eingebaut. Dieser besteht 
aus zwei Elektroden, zwischen diesen 
befindet sich ein flüssigkeitsempfindliches Substrat. Je höher die 
Luftfeuchtigkeit ist, desto mehr Kapazität besitzt dieses Substrat.

Aus diesen Daten ermittelt der Sensor die tatsächlichen Werte durch eine Kalibrierung, 
welche durch den Hersteller vorgenommen wurde.

Er ist in der Lage Temperaturen von 0°C bis 50°C bei einer Messgenauigkeit von  $\pm$ 2,0°C
zu erfassen 
die Luftfeuchtigkeit kann im Bereich von 20\% bis 90\% gemessen werden, 
wobei die Messgenauigkeit $\pm$ 5\% beträgt \citeM{i}{HammLipp.2023}{}{}.

\section{OSI-Referenzmodell}
\section{Netzwerkprotokolle und Referenzmodelle}

% ---------------------------------------------------------
% 1. EINLEITUNG: WAS IST DAS OSI-MODELL?
% ---------------------------------------------------------

Um die Kommunikation in komplexen Netzwerken wie dem in diesem Projekt aufgebauten IoT-System zu standardisieren, wird in der Literatur häufig das \ac{OSI}-Referenzmodell (Open Systems Interconnection) herangezogen. Es wurde in den 1980er Jahren von der Internationalen Organisation für Normung (ISO) entwickelt, um die Interoperabilität zwischen verschiedenen Netzwerksystemen zu gewährleisten. Es besteht aus sieben aufeinander aufbauenden Schichten, welche von Schicht 7 bis 1 beim Senden und von 1 bis 7 beim Empfang durchlaufen werden. Diese sieben Schichten erfüllen jeweils spezifische Aufgaben, die einen reibungslosen und standardisierten Datenaustausch ermöglichen. Dabei ist zu betonen, dass das Modell selbst keine spezifischen Protokolle oder Komponenten definiert, sondern lediglich die Aufgaben der jeweiligen Schichten vorgibt. Aus diesem Grund lassen sich in der Praxis genutzte Protokolle nicht immer eindeutig einer einzigen Schicht zuordnen, da sie oftmals mehrere Aufgaben gleichzeitig übernehmen \citeM{i}{tanenbaum2011}{41}{42}.

% ---------------------------------------------------------
% 2. DIE UNTEREN SCHICHTEN
% ---------------------------------------------------------
\subsection{Physikalische Schicht (Schicht 1)}

In der untersten Schicht, der Physikalischen Schicht, werden die rohen Bits über das Medium übertragen. Es geht also um die elektrotechnische Umsetzung der Datenübertragung im Fall dieses Projekts über WLAN, bzw. die Funkwellen, welche der ESP32 nutzt, um mit dem Netzwerk zu kommunizieren. Aber auch LAN oder USB-Verbindungen würden hier eingeordnet werden \citeM{i}{tanenbaum2011}{43}{43}.

\subsection{Sicherungsschicht (Schicht 2)}

Während die physikalische Schicht lediglich den Bitstrom überträgt, dient die Sicherungsschicht (Data Link Layer) der Abstraktion hardwarenaher Übertragungsfehler. Ziel ist es, der übergeordneten Vermittlungsschicht eine zuverlässige, logische Verbindung zur Verfügung zu stellen. Dieser Prozess basiert primär auf der Framing-Technik, bei der Datenströme in diskrete Einheiten segmentiert werden.

Zur Gewährleistung der Datenintegrität implementiert die Schicht Mechanismen zur Fehlererkennung sowie – je nach Protokollspezifikation – Bestätigungsverfahren (Acknowledgements). Eine integrierte Flusssteuerung schützt zudem langsame Empfängerinstanzen vor einer Datenüberlastung durch den Sender. Im Kontext des ESP32 und der verwendeten WLAN-Kommunikation übernimmt die Sicherungsschicht zudem die Medienzugriffssteuerung (MAC). Diese koordiniert den Zugriff auf das geteilte Funkmedium, um Kollisionen zu minimieren und eine strukturierte Datenübermittlung im Netzwerkverbund sicherzustellen \citeM{i}{tanenbaum2011}{43}{43}.
% ---------------------------------------------------------
% 3. DIE NETZWERKSCHICHT (Wichtig für Docker!)
% ---------------------------------------------------------
\subsection{Vermittlungsschicht (Schicht 3)}
% DEINE AUFGABE: Erkläre, dass es hier um das Routing von Paketen quer durch verschiedene Netzwerke geht. 
% PROJEKTBEZUG: Nenne das IP-Protokoll (Internet Protocol). Erwähne kurz, dass diese Schicht später wichtig wird, um zu verstehen, wie das Docker-Bridge-Netzwerk intern die IP-Adressen für InfluxDB und Grafana routet. \cite{tanenbaum2011}

% ---------------------------------------------------------
% 4. DIE TRANSPORTSCHICHT (Dein wichtigster Teil!)
% ---------------------------------------------------------
\subsection{Transportschicht (Schicht 4)}
% DEINE AUFGABE: Erkläre den Unterschied zwischen TCP (verbindungsorientiert, zuverlässig) und UDP (schnell, unzuverlässig). 
% PROJEKTBEZUG: Erkläre, dass euer System zwingend TCP nutzt, damit keine Sensordaten verloren gehen. Erwähne unbedingt den TCP-Three-Way-Handshake (SYN, SYN-ACK, ACK) - das brauchst du später für die Fehlersuche! \cite{tanenbaum2011}

% ---------------------------------------------------------
% 5. DIE ANWENDUNGSSCHICHT (Zusammenfassung)
% ---------------------------------------------------------
\subsection{Anwendungsschicht (Schicht 5 bis 7)}
% DEINE AUFGABE: Erkläre kurz, dass die Schichten 5 und 6 in der Praxis oft direkt von der Anwendung mitübernommen werden. 
% PROJEKTBEZUG: Hier arbeiten eure eigentlichen Daten-Protokolle: MQTT für das Senden der Sensordaten und HTTP(S) für den Abruf des Grafana-Dashboards.

% ---------------------------------------------------------
% 6. DER SCHWENK ZUM TCP/IP-MODELL
% ---------------------------------------------------------
\subsection{Einordnung in das TCP/IP-Referenzmodell}
Obwohl das OSI-Modell theoretisch präziser ist, hat sich in der Praxis der modernen Netzwerktechnik das TCP/IP-Referenzmodell durchgesetzt \cite{tanenbaum2011}. 
% DEINE AUFGABE: Erkläre in 2-3 Sätzen, dass dieses Modell nur 4 Schichten hat (Netzzugang, Internet, Transport, Anwendung). Es fasst die OSI-Schichten 1+2 sowie 5-7 zusammen. Schließe damit ab, dass die in diesem Projekt genutzte Architektur (IP, TCP, MQTT/HTTP) exakt diesem TCP/IP-Stack entspricht.


\section{MQTT und HTTP}
Der Vergleich!

MQTT: Leichtgewichtig, Publish/Subscribe, perfekt für den ESP32 und instabiles WLAN.
HTTP: Request/Response-Modell, schwergewichtiger. Wird intern von Telegraf (um Daten an InfluxDB zu senden) und von Grafana (für das Web-Dashboard) genutzt.

\section{Container-Netzwerke(Docker Bridge)}
Wie kommunizieren Docker-Container? Erklärung des virtuellen Switches (docker0), der ein komplett eigenes Subnetz aufspannt, getrennt vom physischen WLAN.


\section{Tunneling und Reverse Proxies (Cloudflare Tunnel)}
Die Theorie, wie man ein lokales Netzwerk sicher ins Internet bringt. Erklärung, dass der Tunnel von innen nach außen aufgebaut wird, wodurch man NAT-Firewalls am Router umgeht, ohne Ports öffnen zu müssen.